\section{Analytical and Nonsmooth Benchmarks}\label{sec:benchmarks}
The time-additive Merton model has wealth dynamics
\[
 dX=[(r+p(\mu-r))X-c]dt+p\sigma X\,dW,
 \qquad U(c)=\frac{c^{1-\gamma}}{1-\gamma}.
\]
For the stationary infinite-horizon benchmark, set $c=mX$, $V=KX^{1-\gamma}/(1-\gamma)$, and impose the usual admissibility and transversality requirements. Writing
\[
 p^*=\frac{\mu-r}{\gamma\sigma^2},\qquad
 R_*=r+p^*(\mu-r)-\tfrac12\gamma(p^*)^2\sigma^2,
\]
substitution and maximization give
\begin{equation}\label{eq:merton}
 m^*=\frac{\rho+(\gamma-1)R_*}{\gamma},\qquad K=(m^*)^{-\gamma}.
\end{equation}
With $\rho=.04$, $\gamma=2$, $\mu=.08$, $r=.02$, and $\sigma=.2$, these numbers are $p^*=.75$ and $m^*=.04125$. These are stationary constants; an arbitrary finite-horizon problem with an unspecified bequest does not have the same constant consumption rule. A finite-horizon verification test can use the matching bequest $KX^{1-\gamma}/(1-\gamma)$, explicitly declared rather than silently omitted.

For $\mu=.01$ with no short sales, the unconstrained share is $-.125$ and the constrained optimum is zero at every positive wealth. Its effective return is $r$ and $m^*=.03$. The maximized HJB remains an equality. The action KKT condition is an inequality at the boundary, but it does not turn this regular-control problem into an obstacle problem or create a wealth-dependent kink. The exit experiment in Section~\ref{sec:theory}, rather than this homothetic corner solution, supplies the nonsmooth selection test.

\subsection{Executed calibration and interpretation}
The replication program trains a homothetic neural subfamily: a positive value coefficient, a sigmoid consumption-share output, and a projected portfolio-share output. These are neural output parameters in a known invariant representation, not an unrestricted multilayer approximation. The two Adam optimizers have learning rates $.025$ and $.008$; each outer iteration contains three critic updates and one actor update. There are 1,600 outer iterations and seeds $0,1,2$. All attempted runs, including their coefficient and policy traces, are saved. These calibrations optimize the stationary coefficient equation and then check the resulting coefficient against the matching bequest; the coefficient error is a boundary error in that finite-horizon interpretation, not a boundary condition hard-imposed during these particular updates.

\begin{table}[t]
\centering
\caption{Corrected portfolio calibration: analytical targets and executed block updates}\label{tab:portfolio}
\begin{tabular}{lrrrr}
\toprule
Model & Target $m$ & Target $p$ & Runs & Maximum policy error\\
\midrule
Time-additive Merton & .04125 & .75000 & 3 & $5.6\times10^{-16}$\\
No-short-sale corner & .03000 & .00000 & 3 & $2.3\times10^{-16}$\\
Recursive matching-bequest test & .04550 & .30000 & 3 & $7.0\times10^{-17}$\\
\bottomrule
\end{tabular}
\par\smallskip\footnotesize Policy error is the maximum absolute error in the two constant policy outputs, maximized over all three seeds. Calculations use CPU float64 in the declared homothetic subfamily. Machine-level calibration accuracy is not evidence about an unrestricted network, a GPU runtime, or a general recursive infinite-horizon problem.
\end{table}

Table~\ref{tab:portfolio} reports quantities generated by the new code. The normalized policy-evaluation residual is below $4\times10^{-16}$ in these runs. These results establish that the revised graph and formulas agree in the exactly representable laboratory. They do not establish a broad advantage over a direct analytical calculation. In the exit laboratory, grids with 40, 80, 160, and 320 intervals recover the grid value to within $3\times10^{-15}$; the wrong smooth sequence has scaled residual between 1.984 and 1.999. All grid nodes, not a selected interior sample, enter that diagnostic.

\section{Recursive Utility}\label{sec:recursive}
Let $a=1-1/\psi$, $b=1-\gamma$, and $Y=bV>0$, with $\psi\ne1$ and $\gamma\ne1$. We use the difference-form aggregator
\begin{equation}\label{eq:ez}
 f(c,V)=\frac\rho a\left[c^a(bV)^{1-a/b}-bV\right].
\end{equation}
Its consumption derivative is
\begin{equation}\label{eq:ezderivative}
 f_c=\rho c^{a-1}(bV)^{1-a/b}>0.
\end{equation}
When $\gamma=1/\psi$, one has $a=b$ and $f=\rho(c^b/b-V)$, a positively normalized time-additive CRRA aggregator. The factor $b$ must not multiply the consumption term outside its power. Doing so reverses consumption monotonicity when $\gamma>1$. This distinction follows the difference-form convention discussed by \citet{herdegen2021}.

For a homothetic candidate $V=KX^b/b$, the maximization equation becomes
\begin{equation}\label{eq:eznormalized}
 0=\sup_{m,p}\left\{\frac\rho a[m^aK^{-a/b}-1]
       +r+p(\mu-r)-m-\tfrac12\gamma p^2\sigma^2\right\}.
\end{equation}
The consumption first-order condition is $\rho m^{a-1}K^{-a/b}=1$. Substitution yields
\begin{equation}\label{eq:ezcandidate}
 p^*=\frac{\mu-r}{\gamma\sigma^2},\quad
 m^*=\psi\rho+(1-\psi)\left[r+\frac{(\mu-r)^2}{2\gamma\sigma^2}\right],\quad
 K=[\rho(m^*)^{a-1}]^{b/a}.
\end{equation}
For the retained parameters $\gamma=5$, $\psi=1.5$, $\rho=.04$, $\mu=.08$, $r=.02$, $\sigma=.2$, this gives $(m^*,p^*)=(.0455,.3)$.

\subsection{A well-defined validation problem at the retained parameters}
The numerical calibration uses \eqref{eq:ezcandidate} as a finite-horizon matching-bequest test, not as an unqualified infinite-horizon theorem. Choose $T>0$, stop wealth on leaving $[\underline x,\bar x]$ with $0<\underline x<\bar x<\infty$, and impose the candidate $Kx^b/b$ on the stopping boundary and at $T$. Consumption shares are in a compact positive interval containing $.0455$, and portfolio shares are in a compact interval containing $.3$.

For these retained parameters, $a>0$ and $b<0$. Since $c$ ranges over a compact positive interval, there exist $0<\underline Y<\overline Y<\infty$ such that $f(c,\underline Y/b)\leq0$ and $f(c,\overline Y/b)\geq0$ for every feasible $c$. Indeed, $c^aY^{-a/b}$ tends to zero as $Y\downarrow0$ and to infinity as $Y\uparrow\infty$. Enlarge the interval so that it contains the boundary candidate. The constants $\overline Y/b$ and $\underline Y/b$ are respectively lower and upper barriers for the backward utility equation. Scalar BSDE comparison therefore keeps every policy utility in this interval. The aggregator is Lipschitz there. The smooth homothetic candidate maximizes the Hamiltonian and matches all stopping payoffs; Theorem~\ref{thm:verification} identifies it as the value of this specified finite-horizon problem.

The choice of a matching bequest is an analytical validation device. It does not establish that the same candidate is selected by every infinite-horizon recursive formulation. A separate infinite-horizon application must check its utility-process class and transversality, particularly when risk aversion and intertemporal-complementarity parameters straddle one \citep{herdegen2021}. The paper retains that economic application, but does not use a scalar root to bypass its selection issue.

A critic can enforce $bV>0$ through $b v_\xi=\exp(N_\xi)$. With a positive terminal target, logarithmic interpolation between $bG$ and the neural output enforces both the terminal condition and the power domain. Additional barriers or bounded transformations enforce a declared invariant interval. Such transformations are stated in the architecture, not imposed retrospectively after a failed evaluation.

\subsection{Stability and baseline design}
The recursive experiment in Table~\ref{tab:portfolio} tests the corrected sign, coefficient, and block graph. It does not compare independent DDPG training trajectories. The reviewed manuscript's constructed oscillating curve and manually entered coordinates remain in the archival source, but have no evidentiary role in this revision. A matched neural comparison must use the same utility process, domain, bequest, action restrictions, training budget, and held-out evaluation quantities. It must report all failed seeds and value-domain violations. Neither a zero-centered target nor three successful homothetic calibrations establishes general relative stability.

\section{Endogenous Preferences and Portfolio Choice}\label{sec:ndu}
The endogenous-preference application retains its internal adjustment margin. Let $u$ be risk aversion, $X$ wealth, and $\theta$ adjustment effort. On a declared interval $[\underline u,\bar u]\subset(1,\infty)$, specify the reflected process
\begin{align}
 du&=\theta\,dt+\sigma_u\,dW^u+dL^{\underline u}-dL^{\bar u},\label{eq:ndustate}\\
 dX&=[(r+p(\mu-r))X-c]dt+p\sigma_X X\,dW^X,
 \qquad d\langle W^u,W^X\rangle_t=\varrho\,dt.
\end{align}
The nondecreasing local-time regulators act only at their respective boundaries. This supplies the state mechanism that a bounded adjustment actor does not supply. The numerical reference uses $[1.2,3]$, so the utility singularity at $u=1$ is never reached. In an alternative positive-state model one could specify a multiplicative or transformed diffusion, but that would be a different primitive and must be identified as such.

\subsection{Cardinal utility and the objective}
Fix a reference consumption unit $\bar c>0$. The baseline flow utility is
\begin{equation}\label{eq:nduutility}
 U(c,u)=\frac{(c/\bar c)^{1-u}}{1-u},\qquad
 r_0(c,u,\theta)=U(c,u)-\tfrac12 k\theta^2.
\end{equation}
The reference calculation takes $\bar c=1$, preserving the unnormalized CRRA family of the original application. The agent maximizes
\begin{equation}\label{eq:nduobjective}
 \E_{t,u,X}\left[\int_t^\tau e^{-\rho(s-t)}r_0(c_s,u_s,\theta_s)ds
        +e^{-\rho(\tau-t)}G(u_\tau,X_\tau)\right].
\end{equation}
Both the payoff and the adjustment cost have the same time-relative discount convention. The terminal/exit payoff and the state domain are stated for each calculation.

A useful family for normalization sensitivity is
\[
 U_\kappa(c,u)=\frac{(c/\bar c)^{1-u}-1}{1-u}+\kappa(u).
\]
The baseline corresponds to $\kappa(u)=1/(1-u)$. The log-continuous normalization corresponds to $\kappa=0$ and is an economically different preference-formation model. A calibration of $\kappa$ fixes the cardinal value of a reference consumption stream at each preference state. The choice is not identified by risk-aversion curvature alone.

In particular,
\begin{equation}\label{eq:uderivative}
 U_u(c,u)=\frac{(c/\bar c)^{1-u}\{1-(1-u)\log(c/\bar c)\}}{(1-u)^2}.
\end{equation}
For $c/\bar c>1$ and $u>1$ this derivative is positive. At $c=2$, $u=2$, $\bar c=1$, it is approximately $.84657$. This is a derivative of flow utility, not a sign theorem for the endogenous value derivative $V_u$. The agent's full incentives also depend on future consumption, boundary behavior, and the adjustment technology.

\subsection{Hamiltonian, constrained policies, and economic implications}
The HJB equation is
\begin{align}\label{eq:nduhjb}
 -V_t=\sup_{c,p,\theta}\big\{&U(c,u)-\tfrac12 k\theta^2-\rho V
 +\theta V_u+[(r+p(\mu-r))X-c]V_X\nonumber\\
 &+\tfrac12\sigma_u^2V_{uu}+\tfrac12p^2\sigma_X^2X^2V_{XX}
 +\varrho\sigma_u\sigma_X pX V_{uX}\big\}.
\end{align}
Reflection gives $V_u=0$ at the two preference boundaries, in the appropriate boundary-value sense. The correlation term is retained with its correct coefficient.

For an interior optimum with $V_X>0$ and $V_{XX}<0$, the first-order conditions are
\begin{align}
 c^*&=\bar c\,[\bar c V_X]^{-1/u},\label{eq:nduc}\\
 \theta^*&=V_u/k,\label{eq:ndutheta}\\
 p^*&=-\frac{V_X(\mu-r)}{XV_{XX}\sigma_X^2}
       -\frac{\varrho\sigma_u V_{uX}}{XV_{XX}\sigma_X}.\label{eq:ndupi}
\end{align}
For independent box constraints, strict concavity in the corresponding action permits projection of these unconstrained maximizers onto their intervals. If $V_{XX}$ has the wrong sign or constraints couple the actions, a joint feasible maximization is required instead. The formula for $p$ must not be used where its denominator vanishes.

Equation~\eqref{eq:ndutheta} gives a conditional shadow-price interpretation of internal adjustment. Holding $V_u$ fixed, a larger $k$ reduces its magnitude. Across economies, however, $V_u$ itself depends on $k$, so this is not a global policy comparative static. A valid unconditional implication is that the optimized value is weakly decreasing in $k$: every fixed admissible policy loses $\tfrac12(k_2-k_1)\E\int e^{-\rho s}\theta_s^2ds$ when the cost rises. When an envelope derivative exists,
\begin{equation}\label{eq:envelope}
 \partial_k V^k=-\tfrac12\E^{a^k}\int_t^\tau e^{-\rho(s-t)}\theta_s^2ds\leq0.
\end{equation}
The sign of the hedging component in \eqref{eq:ndupi} remains conditional on $V_{uX}$ and the curvature of wealth value. A negative correlation alone does not determine it.

\subsection{Independent two-state reference and identification}
The executable reference solves the finite-horizon problem with $T=1$, $u\in[1.2,3]$, wealth stopped at $[.5,2]$, and $G(u,X)=X^{1-u}/(1-u)$ at exit and terminal time. Parameters are $(r,\mu,\sigma_X,\sigma_u,\varrho,\rho)=(.02,.08,.2,.08,-.3,.04)$. The finite action set is the product of $c/X\in\{.02,.05,.1\}$, $p\in\{0,.375,.75\}$, and $\theta\in\{-.15,0,.15\}$. A four-point Brownian cubature with positive interpolation weights performs backward induction. The wealth boundary is imposed at every time step; preference overshoots are reflected. This reference defines a finite-action numerical approximation, not continuous-action optimality.

At $(u,X)=(2.1,1)$, the original $17\times25$ state grid and 20 time steps give values $-11.915211$ and $-11.956983$ for $k=1$ and $k=5$. The $25\times37$ grid and 40 time steps give $-11.866768$ and $-11.908738$. Both select $(c/X,p,\theta)=(.1,.75,-.15)$ at that point. The approximately $.048$ refinement change is reported separately from training error. The higher-cost value is weakly lower at every recorded node and time. These original references remain useful consistency checks; Section~\ref{sec:ndu-neural} supplies actual neural fits and shows how strongly the finite action domain affects their economic comparison.

The model also specifies what a structural estimation exercise would require. For observations $Y_j=\mathcal M(S_{t_j},a_{t_j};\vartheta)+\varepsilon_j$, let $m(\vartheta)$ denote selected population moments of the induced observation process. Local identification requires a rank condition on $D_\vartheta m$, after accounting for latent preference states, normalization parameters, and observation noise. Variation of a simulated policy with $k$ alone is insufficient. Numerical value and policy errors must additionally be propagated into moment errors. The endogenous-preference application is therefore retained as an economic model with an explicit internal margin, not used as evidence that its structural parameters have already been identified from data.

The next comparison is the retained R8 evidence under its original enclosure. The new signed envelope and fresh nested-grid policies are reported separately in Section~\ref{sec:r9results}; historical failure counts are not overwritten.
\subsection{Neural preference policies and the role of the action oracle}\label{sec:ndu-neural}
We compute the original preference-adjustment economy with the cardinal utility, shock correlation, reflected preference state, and wealth stopping payoff specified above. The main grid has 17 preference points, 25 log-wealth points, and 20 time steps over a unit horizon. Adjustment cost is $k=1$. The expanded control box is
\[
 m=c/X\in[.02,.30],\qquad p\in[0,1.5],\qquad\theta\in[-.45,.45].
\]
The inherited finite-action comparison uses the 175 triples formed by five consumption shares, five portfolio shares, and seven adjustment efforts, including all endpoints. This is the same constrained economy across the matched finite methods, not a change of utility normalization.

Two experiments distinguish action representation from action optimization. First, the R7 finite study trains two-hidden-layer, width-32 tanh critics and categorical actors. At each backward step it computes the complete 175-action payoff vector using the current continuation critic. Its exact argmax provides a label, and the vector also supplies the actor's expected action-loss objective. Thus \emph{the training oracle is exhaustive}, even when the final actor is accurate. It is a useful neural policy-representation experiment, but not evidence that the actor avoided Bellman maximization.

Second, R8 trains a continuous-output actor directly through its own positive-weight payoff. It never receives an all-action value vector or an exact argmax label. The actor receives 160 Adam updates per time level; the continuation critic receives 320 Adam updates and up to 300 L-BFGS iterations, with failed polishing steps rolled back. The actor-free direct comparator optimizes eight independent bounded controls at every state for the same 160 iterations, followed by the same critic fit. Every restart and rejected action query is charged. Neither procedure invokes a Riccati solution, an analytically invariant policy class, or the independently computed reference during training.

\paragraph*{Finite-action neural policies.}
The raw R7 actor is the primary neural representation result. Its full dynamic policy loss is below the declared $.1$ target for each fixed seed. A top-four shortlist and a separately reported correction map further reduce the loss. The correction map is not part of the raw network. The complete finite policy, not only one-step deviations against a fitted critic, is evaluated by independent backward recursion.
\begin{table}[htbp]\centering
\caption{Finite preference policies: R7 evidence}
\begin{small}\begin{tabular}{lrrrrrr}
\toprule
Method & Seed & Raw loss & Final loss & Value error & Seconds & Guard (\%) \\
\midrule
NBO & 11 & 0.0445 & 0.0261 & 0.0924 & 30.59 & 0.294 \\
Direct & 11 & 0.0181 & 0.0181 & 0.0790 & 24.79 & 0.000 \\
NBO & 29 & 0.0500 & 0.0260 & 0.0814 & 30.49 & 0.332 \\
Direct & 29 & 0.0263 & 0.0263 & 0.0551 & 25.19 & 0.000 \\
NBO & 47 & 0.0500 & 0.0263 & 0.1325 & 29.89 & 0.230 \\
Direct & 47 & 0.0264 & 0.0264 & 0.0824 & 25.17 & 0.000 \\
\bottomrule\end{tabular}\end{small}
\par\smallskip\begin{minipage}{.96\linewidth}\footnotesize Loss bounds cover every stored state and time of the same 175-action economy. Seconds include training and the three reported policy verifications; classical backward induction averages 0.102 seconds. A raw actor, shortlist, and corrected policy are distinct objects. Utility losses are absolute units.\end{minipage}
\end{table}

These results repair the earlier R6 actor failures, which remain in the historical numerical record. They do not reverse the cost comparison: the matched direct finite-Bellman fit is faster and more accurate on average, and classical backward induction is substantially cheaper on this small grid. Tables report those comparisons without selecting a favorable seed. The actor's approximation cost is part of the result, not an argument for discarding the economic application.

\paragraph*{Continuous-action optimization without exact labels.}
The initial continuous actor and direct-gradient pilot fail the $.1$ deviation diagnostic. The actor can saturate at a portfolio bound while a lower portfolio share improves continuation value. Four actor heads with output reinitialization do not cure that failure in the retained pilot. We therefore execute a disclosed local-improvement variant: compare the actor proposal with eight box corners, then perform two coordinate-search sweeps at each of eight decreasing step sizes. The same local search is applied to the direct comparator. An actor-free ablation omits all gradient action updates and starts the identical search from feasible proposals. This ablation measures how much of the hybrid result comes from the search itself.
\begin{table}[htbp]\centering
\caption{Continuous-control training without exact maximizing labels}
\begin{small}\begin{tabular}{lrrrrr}
\toprule
Method & Seed & Raw diagnostic & Hybrid gain & Train (s) & Queries ($10^6$) \\
\midrule
Actor + search & 11 & 0.2596 & 0.0617 & 15.01 & 2.261 \\
Actor + search & 29 & 0.0733 & 0.0188 & 28.49 & 2.261 \\
Actor + search & 47 & 0.2597 & 0.0475 & 26.46 & 2.261 \\
Direct + search & 11 & 0.1554 & 0.0303 & 32.02 & 11.841 \\
Direct + search & 29 & 0.1566 & 0.0203 & 32.80 & 11.841 \\
Direct + search & 47 & 0.1554 & 0.0251 & 30.52 & 11.841 \\
Search only & 11 & 5.5370 & 0.0617 & 13.97 & 0.961 \\
Search only & 29 & 7.2093 & 0.0616 & 13.72 & 0.961 \\
Search only & 47 & 8.0316 & 0.0616 & 12.87 & 0.961 \\
\bottomrule\end{tabular}\end{small}
\par\smallskip\begin{minipage}{.96\linewidth}\footnotesize Raw diagnostic is the largest finite-reference payoff minus the raw policy payoff. Hybrid gain is the augmented finite-deviation gain, a lower bound on continuous-action regret, not a certificate. All methods use zero exhaustive training backups and zero exact argmax training labels. Search-only omits action-gradient updates. Training, reference, and full-cover verification costs are separate.\end{minipage}
\end{table}

The table distinguishes three quantities. The raw-network loss against the finite reference tests the actor alone. The augmented finite-deviation gain permits all 175 inherited actions and the learned continuous action at every state and time; it is a \emph{lower bound} on missed continuous-action gains. The complete action-box upper bound in Proposition~\ref{prop:r8-envelope} is a separate verification result. A small augmented gain is not inserted into that proposition as if it were a uniform upper bound. Local changes occur at nearly every active state-time node in the principal hybrid runs, so this experiment does not describe a sparse correction map or an actor-only controller.
\begin{table}[htbp]\centering
\caption{Full continuous-action policy bounds on the two-state nodal economy}
\begin{small}\begin{tabular}{lrrrrr}
\toprule
Policy & Seed & Uniform bound & Center bound & Eval. (s) & Bound $\leq .1$ \\
\midrule
Actor + search & 11 & 0.1294 & 0.0601 & 0.22 & No \\
Actor + search & 29 & 0.1227 & 0.0583 & 0.22 & No \\
Actor + search & 47 & 0.1301 & 0.0608 & 0.25 & No \\
Direct + search & 11 & 0.1153 & 0.0604 & 0.23 & No \\
Direct + search & 29 & 0.1156 & 0.0586 & 0.22 & No \\
Direct + search & 47 & 0.1203 & 0.0609 & 0.25 & No \\
Search only & 11 & 0.1196 & 0.0647 & 0.24 & No \\
Search only & 29 & 0.1179 & 0.0586 & 0.23 & No \\
Search only & 47 & 0.1209 & 0.0610 & 0.22 & No \\
\bottomrule\end{tabular}\end{small}
\par\smallskip\begin{minipage}{.96\linewidth}\footnotesize The full action box is covered at all 425 states and 20 times. The common optimal upper envelope costs 1624.59 seconds across the retained implementations and 40.96 million action-box evaluations. Each row adds its own interval policy evaluation. Unresolved leaves retain their upper bounds. No continuous-state/time error is set to zero.\end{minipage}
\end{table}

The action cover includes binding constraints and wealth exits. Its unresolved boxes remain in the bound when the declared budget is exhausted. The resulting upper array can audit the actor, the direct optimizer, and the local-search ablation without retraining. This produces an operational error account on a multidimensional state grid with a continuous, nonconcave action objective. It does not establish that all runs meet the tighter $.1$ complete-action target, nor that the actor dominates the direct method at matched economic accuracy.

\paragraph*{Alternative representations and economic units.}
Additional same-model comparisons use a gated neural continuation architecture and tensor Chebyshev projection. The gated method uses DGM-style state-dependent gates \citep{sirignano2018}, but here it is fitted to finite Bellman targets, not called a mesh-free differential DGM solver. The projection exploits the same two-state geometry and imposes the same wealth-boundary lifting. Degrees, parameter counts, all-action costs, failures, and actual accuracy appear in the supplement. These are executed post-review comparisons, not results attributed to the earlier unexecuted development protocol. A generic semilinear BSDE comparator is not substituted for this controlled-covariance, stopped problem without reformulating its second-order control term.

All losses are in the specified cardinal utility units. The optimal finite-model value at the center is approximately $-4.8410$, and the initial value range is approximately $[-9.1599,-.1250]$. A loss of $.1$ is about $2.066$ percent of the center's \emph{absolute utility magnitude}; this scaling is not a consumption-equivalent welfare percentage. Endogenous curvature and adjustment costs make a common consumption-equivalent conversion a separate economic calculation.

\paragraph*{Action bounds and controlled refinements.}
The inherited sensitivity study retains both adjustment-cost values, the original 27-action truncation, 175 and 1,053 actions in the expanded box, a wider action box, wider wealth bounds, and finer space--time grids. Consumption and portfolio upper bounds remain frequently active. Expanding the box changes the constrained economy; refining actions inside one box does not. R8 additionally evaluates the \emph{same frozen continuous policy}, bilinearly interpolated in state and held on its original time schedule, on three nested state--time levels. This isolates changes in policy evaluation from retraining. The supplemental differences are observed refinement diagnostics, not an estimated order promoted to a continuous-time error bound. The original continuous control problem, its utility normalization, and its comparative-static questions remain the economic objects of interest.


\section{Sophisticated Temporal Selves}\label{sec:temporal}
To define present bias in continuous time, fix the discount kernel
\begin{equation}\label{eq:kernel}
 h(s)=\beta e^{-\rho s}+(1-\beta)e^{-(\rho+\kappa)s},
 \qquad 0<\beta\leq1,\quad\kappa>0.
\end{equation}
It satisfies $h(0)=1$ and discounts the immediate future more heavily than the distant future when $\beta<1$. It is a finite-duration continuous-time specification; it is not the assertion that multiplying current utility by $\beta$ creates present bias.

For a candidate Markov policy $a^*$, define the two continuation functions
\begin{equation}\label{eq:components}
 W_j(t,x)=\E_{t,x}^{a^*}\left[\int_t^T e^{-\rho_j(s-t)}u(a_s^*)ds
             +e^{-\rho_j(T-t)}G(S_T)\right],\quad
 (\rho_1,\rho_2)=(\rho,\rho+\kappa).
\end{equation}
Write $V=\beta W_1+(1-\beta)W_2$. Each component evaluates the \emph{same} future policy:
\begin{equation}\label{eq:extended}
 0=(W_j)_t+\LL^{a^*}W_j+u(a^*)-\rho_jW_j,\qquad W_j(T)=G.
\end{equation}
The current self's spike-deviation condition is
\begin{equation}\label{eq:spike}
 a^*(t,x)\in\argmax_{a\in\Act(t,x)}\{u(a)+\LL^aV(t,x)\}.
\end{equation}
The deviations change both current flow utility and the transition of the continuation state. They do not hold all dynamic effects at the equilibrium action. The appendix derives \eqref{eq:spike} from a deviation on $[t,t+\epsilon]$ followed by $a^*$, thereby fixing the equilibrium concept.

A temporal-self NBO uses two policy-evaluation critics for \eqref{eq:extended} and one actor for \eqref{eq:spike}. The current actor does not optimize the sum of future selves' parameter objectives. A target actor can stabilize a numerical iteration, but its agreement with the current actor is not an equilibrium certificate. The relevant diagnostics are component evaluation errors and one-shot deviation gains. The benchmark $\beta=1$ collapses to the usual exponentially discounted problem.

\subsection{An equilibrium calculation with endogenous continuation values}
For deterministic wealth $dX=(rX-c)dt$, CRRA utility with $\gamma>1$, and an infinite horizon, consider $c=mX$. Set $b=1-\gamma$ and $\lambda_j(m)=\rho_j-b(r-m)>0$. Then
\begin{equation}\label{eq:temporalroot}
 W_j(X)=\frac{m^b}{\lambda_j(m)}\frac{X^b}{b},\qquad
 1=m\left[\frac\beta{\lambda_1(m)}+\frac{1-\beta}{\lambda_2(m)}\right].
\end{equation}
The expression on the right is strictly increasing in $m$ on its admissible interval, starts at zero, and diverges at the first pole for $\beta>0$. The equilibrium root is unique in this homothetic class. Its derivative with respect to $\beta$ is negative because $1/\lambda_1>1/\lambda_2$. Thus stronger present bias raises consumption in this specified equilibrium, not merely at a fixed shadow price.

With $r=.02$, $\rho=.04$, $\gamma=2$, and $\kappa=.2$, the roots for $\beta=1,.9,.7,.5$ are respectively $.030000$, $.031372$, $.034606$, and $.038760$. The code independently optimizes finite-duration deviations and records the gain divided by duration as the duration shrinks. For $\beta=.7$, this quantity declines from approximately $4.60\times10^{-5}$ at duration $.1$ to $4.66\times10^{-9}$ at duration $.001$. This is a diagnostic of the derived spike equilibrium. It does not assert that every finite-duration commitment deviation is unprofitable, which would be a different solution concept.

\section{Dynamic Games}\label{sec:games}
The dynamic Cournot application retains firm-specific capital and investment. Firm $i$ has
\[
 dK^i=(I^i-\delta K^i)dt+\sigma_K K^i dW^i,\qquad q_i=AK^i,
\]
and payoff rate
\[
 r_i=P_N(K)q_i-I^i-\tfrac12 b_i(I^i)^2,\qquad
 P_N(K)=\max\left\{0,P_0-\frac1{M_N}\sum_jq_j\right\}.
\]
For each experiment, the market-size normalization $M_N$, capital domain, investment constraints, covariance, terminal/liquidation values, and horizon are part of the primitives. Keeping $M_N=1$ is a fixed-market comparison. Taking $M_N=N$ is a growing-market comparison. They do not have the same competitive-limit interpretation.

Given rivals' Markov policies $a^{-i}$, a player's best-response equation maximizes only its own action. In the stationary additive case,
\begin{equation}\label{eq:game}
 \rho_i V_i=\sup_{a_i\in\Act_i}\{r_i(x,a_i,a^{-i}(x))+\LL^{a_i,a^{-i}}V_i\}.
\end{equation}
A Markov-perfect equilibrium is a feasible profile satisfying these simultaneous unilateral conditions, with the same boundary and admissibility conventions. A zero policy-evaluation residual can hold for any consistently valued profile and therefore does not identify this equilibrium.

\subsection{Player-specific differentiation and deviation bounds}
Player $i$'s critic is trained with every policy frozen. Player $i$'s actor then differentiates only
\begin{equation}\label{eq:gameactor}
 -\E_\nu\HH_i^{a_i(\phi_i),\sg(a^{-i})}
       [\sg(V_i),\sg(DV_i),\sg(D^2V_i)].
\end{equation}
The updated parameter is $\phi_i$, not the full vector of all players' parameters. Rival actions are detached, and player $j$'s loss never contributes to player $i$'s actor gradient. An explicit regression test exercises this graph.

For each player, define the unilateral Hamiltonian gap by maximizing $\HH_i$ over $a_i$ while keeping rivals fixed. Applying Theorem~\ref{thm:verification} to that best-response problem bounds its finite-horizon exploitability by $2e^{L_iT}B_i+A_{L_i}(T)(2e_i+\eta_i)$. In a discounted finite-state game, Theorem~\ref{thm:discrete} gives $(2\delta_i+\epsilon_i)/(1-q_i)$. These are player-specific bounds; a sum of critic losses is not a substitute. They remain valid without uniqueness or convergence of the simultaneous policy-update dynamics, provided the reported profile satisfies the stated errors and best-response comparison hypotheses.

The static Cournot diagnostic makes the distinction visible. With payoffs $q_i(1-q_i-q_j)$, symmetric Nash output is $1/3$ per firm. Joint-profit maximization gives a symmetric allocation $1/4$. At the latter profile, a firm optimally deviates to $3/8$, earning a gain of $1/64$. A joint objective can therefore be perfectly optimized and still fail a unilateral equilibrium test. This arithmetic is included as a mandatory regression check.

\subsection{Symmetry and market scaling}
Identical initial weights do not prove equilibrium uniqueness or permutation equivariance. An equivariant actor can be specified using own capital and a permutation-invariant aggregation of rival states. Even then, one must check unilateral deviations rather than infer equilibrium from architectural symmetry. Multiple starting profiles and unsuccessful best-response searches belong in a numerical report.

Under fixed market size, positive constant capital per firm cannot be called the competitive limit while total output tends to infinity and the truncated price tends to zero. Maintaining that capital entails investment and adjustment costs. Under $M_N=N$, positive capital per firm may instead be compatible with a growing market, but it describes a different sequence of economies. The original unsupported capital table remains in the archive. Section~\ref{sec:dynamic-computation} replaces its evidentiary role with a fully specified finite dynamic game and independent best-response calculations, without asserting convergence of unrestricted neural game training.

\subsection{Neural dynamic capital policies and full unilateral deviations}\label{sec:dynamic-computation}
For two firms, we retain the capital game with $A=1$, $P_0=2$, $\delta=.1$, $\sigma_K=.15$, $\rho=.04$, $T=3$, investment in $[0,.8]$, and capital in $[.05,1.5]^2$. Terminal or joint-domain exit pays $.5K_i$. Baseline adjustment coefficients are $b_1=b_2=1$; the earlier asymmetric experiment has $b_2=1.2$. The principal neural grid has $17\times17$ states, 30 time steps, and 11 actions per player. Positive cubature and interpolation define the finite game. Market normalizations $M_2=1$ and $M_2=2$ are distinct economies, not evidence for an uncomputed large-population limit.

At each time and state, the R7 calculation trains separate critics and actors. It constructs both full $11\times11$ continuation-payoff matrices. The training labels are a pure one-step equilibrium when one exists and a minimum-maximum-defect pair otherwise. The actor's own objective is differentiated with the rival action and continuation critic fixed. Thus this study actually trains neural strategic policies, but it also pays for an exhaustive joint-action training oracle. Raw policies, three-by-three shortlist policies, and any corrected policies are distinct stored arrays.

After freezing the entire profile, an independent backward optimization permits every own action at every state and time against the fixed rival policy. Its value minus the profile payoff is the unilateral exploitability map. This is not a one-period test with the approximate continuation held fixed. The raw and final profiles meet the $.1$ absolute-profit target in every principal R7 case; the full-oracle guard is unused in those cases.
\begin{table}[htbp]\centering
\caption{R7 neural dynamic-game profiles and full best responses}
\begin{small}\begin{tabular}{lrrrrrr}
\toprule
Market & Seed & Raw max & Firm 1 & Firm 2 & Missing & Seconds \\
\midrule
1 & 11 & 0.0259 & 0.0164 & 0.0252 & 3 & 52.43 \\
2 & 11 & 0.0241 & 0.0255 & 0.0163 & 1 & 56.13 \\
1 & 29 & 0.0300 & 0.0288 & 0.0296 & 0 & 52.44 \\
2 & 29 & 0.0626 & 0.0624 & 0.0547 & 1 & 57.23 \\
1 & 47 & 0.0190 & 0.0081 & 0.0183 & 0 & 52.53 \\
2 & 47 & 0.0101 & 0.0074 & 0.0113 & 0 & 56.98 \\
\bottomrule\end{tabular}\end{small}
\par\smallskip\begin{minipage}{.96\linewidth}\footnotesize All bounds cover every stored finite state, time, and unilateral action, including deviations at the reported missing pure stages. The full-oracle guard is unused in these principal runs. Profit losses are absolute units; market normalizations are different economies.\end{minipage}
\end{table}

Five state-time nodes across the six neural runs lack a pure one-step equilibrium for the fitted continuation matrices. Their locations and defects are reported in the supplement. A minimum-defect selection at those nodes is not called an exact equilibrium. Nevertheless, the complete dynamic best-response bound remains valid for the resulting profile: its proof does not assume that every fitted one-step game has a pure solution. Moreover, for fixed Markov rivals, randomized or history-dependent own deviations cannot improve on the finite Markov best-response value, because each conditional action mixture is bounded by its largest pure-action payoff. The calculation therefore checks such unilateral deviations as well; it does not compute a mixed-strategy equilibrium profile or a continuous-action game.

R8 adds actor-free neural continuation fitting and classical backward selection for both market normalizations. The neural comparator allocates 480 critic Adam updates per time level rather than R7's 320 critic plus 160 actor updates, with the same L-BFGS cap. Counts and timings reveal the difference between nominal update matching and matched wall time. Both methods receive their own independent full dynamic best responses. The earlier six classical game cases, asymmetric costs, alternative tie rules, boundary frequencies, coarse counterexample, and finer-grid results remain in the supplement and raw record.
\begin{table}[htbp]\centering
\caption{R8 actor-free dynamic-game comparisons}
\begin{small}\begin{tabular}{lrrrrrr}
\toprule
Method & Market & Seed & Firm 1 & Firm 2 & Train (s) & Verify (s) \\
\midrule
Classical & 1 & 11 & 0.0000 & 0.0000 & 0.22 & 6.40 \\
Classical & 2 & 11 & 0.0000 & 0.0000 & 0.23 & 6.29 \\
Direct & 1 & 11 & 0.0223 & 0.0273 & 37.43 & 6.09 \\
Direct & 1 & 29 & 0.0199 & 0.0282 & 37.68 & 6.18 \\
Direct & 1 & 47 & 0.0223 & 0.0409 & 36.03 & 6.26 \\
Direct & 2 & 11 & 0.0091 & 0.0694 & 40.43 & 5.75 \\
Direct & 2 & 29 & 0.0306 & 0.0185 & 43.07 & 6.34 \\
Direct & 2 & 47 & 0.0206 & 0.0358 & 40.47 & 6.08 \\
\bottomrule\end{tabular}\end{small}
\par\smallskip\begin{minipage}{.96\linewidth}\footnotesize The direct method fits two continuation networks with the same class and 480 critic Adam updates, versus R7's 320 critic and 160 actor updates. Classical selection uses no neural approximation. Every method receives a new full dynamic best response. R7 and R8 timings were measured in their recorded environments; cross-environment ratios are not a hardware-controlled speed experiment.\end{minipage}
\end{table}

The finite-game result is a numerically evaluated approximate equilibrium with a specified profit-unit bound. It is not inferred from stationarity of a summed actor loss or the static $1/64$ diagnostic. State interpolation, continuous investment, diffusion timing, and continuously monitored exits remain separate terms in a continuous-game error account.


\section{Scalability and the Diffusion Trace}\label{sec:scaling}
The original multi-stock consumption-savings specification has additive utility, independent transitions, and no common resource constraint. It decomposes as $V(k)=\sum_jv_j(k_j)$. A separability-aware comparator solves the one-dimensional components, and identical components can be reused. An unreduced sparse grid is not an appropriate sole comparator for that economy.

To test a coupled generator, the new laboratory instead uses
\begin{equation}\label{eq:lq}
 dX=(AX+a)dt+\Sigma dW,\qquad
 \sup_a\E\int_0^\infty e^{-\rho t}[-X^\top QX-a^\top a]dt.
\end{equation}
The drift is a nonnormal dense matrix $A=-.7I+.06Z/\sqrt d$, where $Z$ is generated by the fixed model seed $991+d$. The payoff matrix is $Q=\operatorname{diag}(1,\ldots,2)+.3\bm1\bm1^\top/d$. The code verifies that $AQ\ne QA$, so the common independent-coordinate reduction of the original benchmark is unavailable. All methods still exploit the declared quadratic structure. A Riccati solver is therefore the appropriate strong reference, rather than a deliberately unreduced general grid.

For $V=-x^\top Px-c_0$ and $a=-Kx$, policy evaluation solves
\begin{equation}\label{eq:lyapunov}
 (A-K-\rho I/2)^\top P+P(A-K-\rho I/2)=-(Q+K^\top K),
 \quad c_0=\tr(\Sigma\Sigma^\top P)/\rho.
\end{equation}
Hamiltonian improvement gives $K=P$. A positive semidefinite $P$ is representable by a signed sum of squared linear neural units; the actor is linear. The exact block laboratory uses a Lyapunov solve for this critic and exact actor improvement. It is an exactly representable NBO subfamily and a matrix-identity test, not evidence that generic stochastic neural optimization has the same convergence rate.

Across dimensions $2,5,10,20,50$ and three stabilizing actor initializations, the block iteration takes six steps. The maximum error in $P$, in operator norm, is below $10^{-13}$ relative to an independent algebraic Riccati solver. The policy-evaluation coefficient residual is below $10^{-12}$. On the unit Euclidean ball, the actor gap is bounded by $\|K-P\|_2^2$, and the residual bound follows from the operator norm of the quadratic coefficient error. The program saves matrices, iteration traces, and CPU wall times for both methods. These very short, single-process timings are diagnostics and do not establish a fixed-accuracy runtime ranking.

\subsection{Trained, coupled nonquadratic capital economies}\label{sec:coupled-neural}
To separate neural scaling from exact quadratic algebra, consider $d$ positive capital stocks and write $Y_i=\log X_i$. Consumption is $c_i=m_iX_i$, with $m_i\in[.02,2]$, and
\begin{align}\label{eq:coupled-r6}
 dY_i&=\left[.10+.15\tanh((B_dY)_i)-m_i-\tfrac12(.15^2+.10^2)\right]dt
       +.15\,dW_i+.10\,dW_0,\\
 r_0(Y,m)&=d^{-1}\sum_i(\log m_i+Y_i)-.10\bar m^2,
 \qquad \bar m=d^{-1}\sum_i m_i.\nonumber
\end{align}
The dense row-normalized matrix $B_d$ is generated once from a recorded seed. Horizon is one, discount is $.04$, and terminal utility is $\bar Y-.03d^{-1}\sum_i(Y_i-\bar Y)^2$. The actual state domain is $\R^d$; no artificial lateral stopping condition is imposed. Dense nonlinear productivity, shared adjustment cost, common noise, and terminal dispersion couple the coordinates. Neither the solver nor its comparator calls a Riccati or Lyapunov routine, and no exact solution is built into the network.

Strict concavity gives a global Hamiltonian maximizer at each queried state. For $p_i=\partial_i v$, its components satisfy
\[
 m_i^*=\operatorname{clip}\left((dp_i+.2\bar m^*)^{-1},.02,2\right),
\]
with the upper bound selected when the denominator is nonpositive. The scalar consistency equation for $\bar m^*$ is monotone and is solved by bisection. This supplies the same feasible maximization routine to the direct-HJB comparator and to action-gap measurement; it does not supply the value function.

Both networks have two hidden tanh layers of width thirty-two. The critic enforces the terminal payoff exactly. The main design uses three seeds at $d=2,5,10,20$, 600 Adam updates, and batches of 128 independently sampled state--time points. Exact automatic differentiation supplies the trace. The held-out set contains 512 fixed independent points in $[0,1]\times[-.5,.5]^d$. A common diagnostic target is residual root mean square at most $.025$ and maximum queried action gap at most $.01$. The recorded first crossing is not necessarily sustained; final failures are also reported.
\begin{table}[t]\centering
\caption{Actual training in dense nonquadratic capital models}\label{tab:coupled-r6}
\small\begin{tabular}{lrrrrr}
\toprule
Method & $d$ & Worst RMS & Worst gap & Failures & Median seconds\\\midrule
NBO & 2 & 0.0161 & $2.42\times10^{-4}$ & 0/3 & 2.4911\\
Direct HJB & 2 & 0.0161 & 0 & 0/3 & 2.9431\\
NBO & 5 & 0.0131 & 0.0020 & 0/3 & 3.6224\\
Direct HJB & 5 & 0.0129 & 0 & 0/3 & 4.6135\\
NBO & 10 & 0.0188 & 0.0052 & 0/3 & 5.9657\\
Direct HJB & 10 & 0.0168 & 0 & 0/3 & 6.6025\\
NBO & 20 & 0.0320 & 0.0114 & 2/3 & 11.5921\\
Direct HJB & 20 & 0.0170 & 0 & 0/3 & 11.8949\\
\bottomrule\end{tabular}
\par\smallskip\footnotesize Maxima are across three seeds, each evaluated at 512 independent points in $[0,1]\times[-.5,.5]^d$. The failure criterion is RMS residual above $.025$ or queried action gap above $.01$ after 600 updates. Both methods use the same critic architecture and global action routine. Times include training diagnostics, not the subsequent Monte Carlo evaluations. These are sampled diagnostics, not uniform economic bounds.
\end{table}

The experiment measures real neural approximation and optimization costs rather than exact quadratic identities. It also shows why a diagnostic target is not a welfare certificate. Simulated policy paths can leave the fitted state box; Euler policy-payoff estimates from the initial state need not agree with a critic that has a small sampled residual inside that box. A follow-up expands the fitting and verification box to $[-1.5,.5]^d$ and uses 1,000 updates at dimensions ten and twenty. Both designs, paired policy-payoff samples, and forty/eighty-step simulation outputs are retained. Sampling standard errors do not remove time-discretization bias or certify an optimum. Process high-water memory is labeled as such, not interpreted as an isolated per-network allocation. No high-dimensional uniform residual bound is claimed for these sampled quantities.

\subsection{Method-paired payoff evidence and arithmetic replay}\label{sec:r8-coverage}
The R7 follow-up evaluates the coupled-capital policies at dimensions ten and twenty, for all three fixed seeds, both fitting domains, and 40, 80, and 160 Euler steps. The same 512 Brownian paths are used by NBO and the direct comparator within every cell. For path $j$, define $D_j=J_j^{\mathrm{NBO}}-J_j^{\mathrm{direct}}$. We report $\overline D$ and the pointwise Student interval $\overline D\pm t_{511,.975}s_D/\sqrt{512}$. Pairing removes the irrelevant component of common path variation from the comparison. The complete set of 36 comparisons is reproduced from the pinned raw arrays in the supplement.
\begin{table}[htbp]\centering
\caption{Paired payoff differences at the finest recorded Euler level}
\begin{small}\begin{tabular}{lrrrrr}
\toprule
Dim. & Seed & Domain & Difference & 95\% lower & 95\% upper \\
\midrule
10 & 11 & Narrow & -0.03338 & -0.03544 & -0.03133 \\
10 & 11 & Wide & 0.00052 & 0.00049 & 0.00054 \\
10 & 29 & Narrow & -0.03094 & -0.03250 & -0.02938 \\
10 & 29 & Wide & -0.00003 & -0.00004 & -0.00002 \\
10 & 47 & Narrow & -0.02755 & -0.02941 & -0.02569 \\
10 & 47 & Wide & 0.00000 & -0.00001 & 0.00001 \\
20 & 11 & Narrow & -0.06302 & -0.06572 & -0.06033 \\
20 & 11 & Wide & 0.00084 & 0.00080 & 0.00088 \\
20 & 29 & Narrow & -0.01449 & -0.01490 & -0.01408 \\
20 & 29 & Wide & 0.00409 & 0.00394 & 0.00424 \\
20 & 47 & Narrow & -0.20161 & -0.20211 & -0.20111 \\
20 & 47 & Wide & 0.00058 & 0.00053 & 0.00062 \\
\bottomrule\end{tabular}\end{small}
\par\smallskip\begin{minipage}{.96\linewidth}\footnotesize NBO minus direct payoff on the same 512 paths. Full 40/80/160-step comparisons are in the supplement. Confidence intervals quantify Monte Carlo variation only.\end{minipage}
\end{table}

The old twenty-dimensional fitting box is exited by essentially all recorded paths. Wider-domain fitting eliminates observed exits in the recorded audits and changes both payoff agreement and method differences. Zero observed exits in 512 paths is not a zero exit-probability theorem. Likewise, these intervals quantify pointwise Monte Carlo uncertainty under a specified Euler scheme. They are neither simultaneous confidence statements across all 36 comparisons nor bounds on time bias, optimal value, or the global differential residual. The experiment identifies a consequential coverage failure and its measured correction without redefining the unbounded economic state space as a stopped box.

The one-state consumption certificate has a different evidentiary status. R7 independently replays its previously generated complete leaf cover with \texttt{mpmath.iv} at 40 decimal digits, without importing the original interval arithmetic implementation. For the stored seed-11 networks, the NBO policy-loss bound is about $.04961$ and the direct-HJB bound about $.04967$. Both meet the stated threshold. This is an \emph{independent interval-arithmetic replay on the previously generated complete cover}, not independent network training, independent cover generation, or a formal machine proof. Together with the new complete three-control cover on the two-state preference grid, it demonstrates two different ways of operationalizing verification. Neither removes the dimensional cost of covering an unrestricted high-dimensional continuous state space.


\subsection{Exact and stochastic trace computation}
Let $\Sigma$ include any covariance factor and set $B=\Sigma^\top D^2v\Sigma$. For independent probes with identity covariance,
\begin{equation}\label{eq:trace}
 q=\tr(B),\qquad \widehat q=\frac1K\sum_{k=1}^Kz_k^\top Bz_k.
\end{equation}
Gaussian probes have variance $2\|B\|_F^2/K$. If $r=d-q/2$, then
\begin{equation}\label{eq:tracebias}
 \E(d-\widehat q/2)^2=r^2+\tfrac14\operatorname{Var}(\widehat q).
\end{equation}
The extra term depends on the critic. Unbiased trace estimation therefore does not give an unbiased squared-residual objective.

Two conditionally independent probe banks repair this particular bias:
\begin{equation}\label{eq:crossprobe}
 \widehat L=(d-\widehat q^{(1)}/2)(d-\widehat q^{(2)}/2),
 \qquad\E[\widehat L\mid t,x,\xi]=r^2.
\end{equation}
With dominated differentiation, its gradient is also unbiased. The sample product need not be nonnegative and can have substantial variance, so this is an objective identity rather than a guarantee of stable optimization. Exact low-dimensional derivatives remain the independent check. If the actor or deterministic term depends on the same probes, conditional independence must be re-established; the simple identity cannot be reused automatically.

The relevant computational cost depends on the network evaluation and differentiation cost, Brownian rank, covariance factorization, sample count, optimization iterations, and requested accuracy. Forming a dense Hessian involves quadratic storage; Hessian-vector products can avoid it. Neither operation count by itself is a total-complexity theorem. Active control bounds do not universally create gradient discontinuities, and a noisy trace cannot make a correctly projected box action leave its box. Feasibility and value-derivative accuracy are separate diagnostics.

The independent-bank estimator is also exercised during actual training in the coupled ten-state economy. One implementation uses two independent banks of two Gaussian Hessian-vector probes; the matched single-bank ablation squares a four-probe estimate. Exact traces are used only for the common held-out diagnostics. All three seeds, objective signs, optimization histories, and policy-payoff samples are saved. The two training procedures perform similarly at the retained covariance scale; unbiasedness alone does not establish a finite-budget advantage. A separate Monte Carlo test checks both the product objective and its parameter gradient against exact trace expressions with declared tolerances. These calculations validate the estimator's execution without identifying a random training objective with a deterministic certificate.


\section{Validated Evaluation and Action Improvement}\label{sec:r9method}
The neural Bellman map must account for two distinct approximation decisions: which continuation values are passed to the next stage, and how tightly the admissible action supremum is bounded. This section supplies an explicit output rule for the first decision and a signed enclosure for the second. Neither construction presumes that a stationary point of the training loss is economically optimal. Both apply to the original preference economy and leave its controls, state law, and stopping convention unchanged.

\subsection{An evaluation rule with a prescribed nodal error}
Let $T_h^a$ denote a positive-weight, one-step policy operator on a finite state grid, and suppose its continuation weights have sum at most $q_h\leq1$. At a given date, let $w$ be the continuation vector, $a$ the selected feasible policy, $\widehat v$ the fitted neural critic, and $\widetilde T_h^a w$ the computed selected-policy target. We distinguish the computation of this target from maximization over actions. In particular, it uses one action at each node, not a table of optimizing labels. For a prescribed tolerance $\tau_h>0$, define
\begin{equation}\label{eq:r9guard}
 C_h(\widehat v,a,w)_i=
 \begin{cases}
  (\widetilde T_h^a w)_i,& |\widehat v_i-(\widetilde T_h^a w)_i|>\tau_h,\\
  \widehat v_i,&\text{otherwise}.
 \end{cases}
\end{equation}
The comparison in the implementation uses an outward upper bound for the difference. A selected entry is assigned the target directly; it is not reconstructed by adding two cancelling floating-point numbers. The indexed corrections, their masks, and the uncorrected neural outputs are saved. The deployed continuation is therefore a neural representation with a recorded nodal correction, not a raw network.

\begin{proposition}[Evaluation guard]\label{prop:r9guard}
Suppose $\|\widetilde T_h^a w-T_h^a w\|_\infty\leq\eta_h$. Then
\[
 \|C_h(\widehat v,a,w)-T_h^a w\|_\infty\leq\tau_h+\eta_h.
\]
For a frozen sequence of policies over $N$ stages, with exact terminal data, the guarded backward values and the exact values of those policies differ by at most $\sum_{j=0}^{N-1}q_h^j(\tau_h+\eta_h)$.
If $N h=T$, $\tau_h=0.1h^2$, and $\eta_h=o(h)$, this evaluation perturbation vanishes as $h\downarrow0$. This statement is a property of the output rule, not a rate for the neural optimizer or an assertion of policy optimality.
\end{proposition}

The guard gives an implementable meaning to the evaluation component of a vanishing-perturbation condition. It does not supply the action-improvement error, the consistency error of a stopped diffusion approximation, or convergence of an uncorrected network away from the nodes. Its cost can be substantial: in the numerical economy it corrects a large fraction of fitted nodal values. We consequently count the correction memory and compare against an actor-free method with the identical guard. The construction is not evidence of neural compression when almost every node is retained.

\subsection{Signed enclosures for continuous actions}
Fix a state-time node and a continuation vector $w$. Write $Q_w(a)$ for its one-step payoff. The action set is the full compact box of the preference economy. Let $B=[\ell,r]$ be an action rectangle with centre $c$. A natural interval calculation supplies an upper bound $U^{\mathrm{nat}}(B)$, while independent feasible evaluations supply a lower incumbent $L$.

The continuation is bilinear within state cells and continuous across their internal faces. Reflection changes its one-sided slopes. A stopped transition is different: the implementation switches between interpolated continuation and the specified stopping payoff, which need not agree away from boundary grid points. We call $B$ \emph{regular} when every cubature transition stays entirely inside, or entirely outside, each wealth stopping face. Only on such a box do we use derivative information. All intersected interpolation cells and both sides of reflection folds enter the derivative hull.

Suppose $G_j(B)=[g_j^-,g_j^+]$ encloses every one-sided partial derivative with respect to action component $j$ on a regular box. The mean-value upper enclosure is
\begin{equation}\label{eq:r9mean}
 U^{\mathrm{mv}}(B)=\overline Q_w(c)+
 \sum_j\max\{g_j^-(\ell_j-c_j),g_j^-(r_j-c_j),g_j^+(\ell_j-c_j),g_j^+(r_j-c_j)\},
\end{equation}
where $\overline Q_w(c)$ is an outward point upper bound. The algorithm uses the smaller of this bound and $U^{\mathrm{nat}}(B)$. If $g_j^->0$, it restricts the supremum to $a_j=r_j$; if $g_j^+<0$, it restricts it to $a_j=\ell_j$. These are verified reductions to a maximizing face, not decisions based on a sampled derivative. For boxes crossing a stopping switch, the natural union enclosure is retained without this reduction.

\begin{theorem}[Signed action cover]\label{thm:r9cover}
Assume that the natural enclosures contain their respective payoffs, that the derivative hulls contain all one-sided derivatives on each regular box, and that feasible point values and arithmetic bounds are outward. Start from a cover of the action box and apply subdivision, reduction to a verified maximizing face, and removal of a box only when its upper bound is at most $L+\varepsilon$. At any interruption, retain the upper bound of every unresolved box. Then
\begin{equation}\label{eq:r9coverbound}
 \sup_{a\in\Act}Q_w(a)\leq\max\{L+\varepsilon,\ \max_{B\ \mathrm{unresolved}}U(B)\}.
\end{equation}
The bound remains valid when a depth, work, or time budget is exhausted. If $r$ action coordinates have verified strict derivative signs throughout a regular parent box, maximization over that box reduces to a face of dimension at most $d_a-r$.
\end{theorem}

The last statement describes a source of computational improvement, not a dimension-free worst-case bound. Under a $K$-Lipschitz gradient, a regular stationary box of diameter $s$ has mean-value excess of order $Ks^2$ when its point evaluation error is negligible. Boxes that cross stopping switches, have ambiguous derivative signs, or contain many interpolation cells can remain expensive. Our experiment measures their actual cost and keeps unresolved bounds; it does not assume that every box has favorable structure.

The operational NBO output is now explicit: the frozen policy representation, its guarded continuation vector when used, an independently propagated policy interval, and a complete-action optimal upper envelope. A nodal lookup policy is the object certified in the preference experiment. An off-grid interpolation of that policy is a different deployment rule and is evaluated separately. No online search is charged to a stored lookup table; constructing that table includes every actor, critic, and local-search query.

\section{Cost, Refinement, and Economic Accuracy}\label{sec:r9results}
The preceding applications define the original economic scope of NBO. We now evaluate the guarded map and signed cover on the same preference economy, and connect the high-dimensional experiment to an absolute economic benchmark. The comparison retains every earlier failure. In particular, a new upper envelope does not turn an actor-plus-search policy into an actor-only policy, and retraining on a finer grid does not repair the off-grid performance of a previously frozen policy.

\subsection{A tighter certificate for identical frozen policies}
The reference specification has $17\times25$ state nodes, twenty dates, and the continuous action box $[0.02,0.30]\times[0,1.5]\times[-0.45,0.45]$. The ratios constrain consumption, leverage, and preference adjustment in the declared economy. They are economic primitives, not bounds on an unconstrained optimum. Active consumption or portfolio limits remain meaningful features of this constrained specification; changing them changes the economy.

We freeze the nine previously computed actor-plus-search, direct-plus-search, and search-only policies, including their original seeds. The complete-action upper envelope is recomputed independently of these policies' lower payoff envelopes. The new cover uses a one-step tolerance of $0.0002$ and retains every unresolved upper bound. Its output can be shared across policies only because their action set, transition rule, domain, and payoff normalization coincide.
\begin{table}[htbp]
\centering\small\setlength{\tabcolsep}{4pt}
\caption{Complete-action bounds for the same frozen policies}\label{tab:r9common}
\begin{tabular}{lrrrr}
\toprule
Policy & Seed & Uniform bound & Centre bound & $\leq0.1$ \\
\midrule
Actor + search & 11 & 0.0687 & 0.0068 & Yes \\
Actor + search & 29 & 0.0553 & 0.0050 & Yes \\
Actor + search & 47 & 0.0552 & 0.0075 & Yes \\
Direct + search & 11 & 0.0465 & 0.0071 & Yes \\
Direct + search & 29 & 0.0536 & 0.0052 & Yes \\
Direct + search & 47 & 0.0466 & 0.0075 & Yes \\
Search only & 11 & 0.0633 & 0.0113 & Yes \\
Search only & 29 & 0.0655 & 0.0053 & Yes \\
Search only & 47 & 0.0631 & 0.0076 & Yes \\
\bottomrule
\end{tabular}
\par\vspace{3pt}\begin{minipage}{0.97\linewidth}\footnotesize Full $17\times25$ state grid, twenty dates, unchanged three-control box. The common signed cover took 372.53 seconds; 20 of twenty one-step covers closed at the requested tolerance. Other levels retain all unresolved upper bounds. Bounds are conditional on the stated outward-arithmetic execution model; continuous-state/time transfer is not included.\end{minipage}
\end{table}

Table~\ref{tab:r9common} reports the resulting all-date, all-node loss bounds. The improvement relative to the historical table concerns the tightness of the verification calculation, not additional training or a change in preferences. The raw actors remain separately reported in the machine-readable common accounts. Search remains consequential, and these results do not establish an actor advantage over the actor-free comparator. They do show that the earlier failure of a $0.1$ upper-bound target need not imply that the frozen hybrid policies themselves are more than $0.1$ from the nodal optimum.

For a controlled cost comparison, Table~\ref{tab:r9matched} gives one-step corner-only and signed covers at three dates, with identical continuation arrays, initial feasible actions, tolerance, depth limit, and work budget. Wall time includes derivative hull construction for the signed cover. The point of this experiment is not to compare two historical machines. Historical training costs are retained in the supplementary total-cost accounts with that qualification. Shared verification is amortized under explicitly stated populations of one, nine, and one hundred policies, and the welfare conversion has a separately counted cost. No claim of neural speed advantage is inferred from this accounting.
\begin{table}[htbp]
\centering\small\setlength{\tabcolsep}{4pt}
\caption{Matched one-step action-cover computation}\label{tab:r9matched}
\begin{tabular}{lrrrrr}
\toprule
Cover & Date & Seconds & $Q$ boxes, m. & Bracket & Budget hit \\
\midrule
Corner & 19 & 18.67 & 0.414 & 0.01336 & Yes \\
Signed & 19 & 9.50 & 0.066 & 0.00020 & No \\
Corner & 10 & 19.89 & 0.440 & 0.02230 & Yes \\
Signed & 10 & 14.60 & 0.112 & 0.00020 & No \\
Corner & 0 & 18.32 & 0.402 & 0.03430 & Yes \\
Signed & 0 & 25.84 & 0.446 & 0.00173 & Yes \\
\bottomrule
\end{tabular}
\par\vspace{3pt}\begin{minipage}{0.97\linewidth}\footnotesize Continuation, incumbent, tolerance, and budget are identical within each date. Wall time includes gradient work; a query count is not a hardware-independent speed measure.\end{minipage}
\end{table}


\subsection{The whole-domain refinement failure and its repair}
A frozen policy has to be distinguished from a policy refitted on a new grid. The complete historical refinement arrays reveal large losses away from the centre. Their maximum positive reference advantages, locations, quantiles, boundary-region summaries, and action saturation rates are now included in the supplement. The largest historical discrepancy is approximately $2.1234$ at high risk aversion and wealth close to the lower stopping boundary. Centre values alone do not summarize that experiment. These discrepancies are feasible-reference lower bounds on policy regret, not upper bounds for the continuous economy.

The new refinement comparison uses genuinely nested spatial grids, $17\times25$ and $33\times49$, with twenty and forty dates, respectively. Each declared actor and actor-free search case is trained afresh. The network width, activation, actor budget, critic budget, and seed set are unchanged across the relevant ablation. We compare the unguarded fitted critic with the indexed guard in \eqref{eq:r9guard}. For the forty-date case the target tolerance is $0.0000625$. Both the fraction of corrected nodes and their storage cost are reported, rather than being attributed to an uncorrected neural approximation.
\begin{table}[htbp]
\centering\small\setlength{\tabcolsep}{4pt}
\caption{Fresh policies on the nested fine grid}\label{tab:r9nested}
\begin{tabular}{lrrrr}
\toprule
Method & Uniform bound range & Ref. advantage & Mean train s. & Passes \\
\midrule
Actor & 0.0654--0.0963 & 0.0932 & 88.32 & 3/3 \\
Actor + guard & 0.0140--0.0356 & 0.0155 & 87.75 & 3/3 \\
Search & 0.0655--0.0963 & 0.0932 & 51.07 & 3/3 \\
Search + guard & 0.0349--0.0357 & 0.0156 & 47.40 & 3/3 \\
\bottomrule
\end{tabular}
\par\vspace{3pt}\begin{minipage}{0.97\linewidth}\footnotesize Full $33\times49$ state grid, forty dates; seeds 11, 29, and 47. Reference advantage is a lower regret diagnostic against the feasible 175-action reference, not an upper bound. Uniform bounds use a separately computed complete-action fine-grid envelope. The guard is a saved indexed correction, not an uncorrected network.\end{minipage}
\end{table}

A new complete-action envelope is computed on the fine grid, not interpolated from the coarse certificate. It is applied to all twelve fine-grid policies, including the six unguarded runs. This separates three facts: the failure of a frozen coarse-grid policy; the accuracy achieved by retraining within a declared nodal economy; and the additional error required to transfer a result to a continuously monitored diffusion. The first remains in the record, the second is measured, and the third is not assigned a zero value. Two levels alone do not establish an asymptotic convergence rate.

The deployment specification is a stored state-time action table with its saved neural proposal and continuation correction. Evaluating this table at its nodes requires no online maximizing search. Interpolating the table at off-grid states invokes a different policy rule; the historical experiment shows why this distinction matters. A raw actor without a table or local improvement has a different memory requirement and is not represented by the hybrid loss bounds.

\subsection{A model-specific welfare account}
An absolute utility gap is not a consumption-equivalent percentage when the utility exponent is itself a controlled state. We therefore compute a defined compensating transfer. Hold the saved policy and its state law fixed, leave preference-adjustment costs and the stopping bequest unchanged, and add externally financed consumption so that flow consumption becomes $(1+e)c_tX_t$. This intervention changes felicity but does not charge its financing to wealth. It is a compensation experiment, not another admissible self-financed policy in the original box.

Let $J_e^\pi$ be the compensated payoff and $U$ the common optimal upper envelope. A verified lower envelope $\underline J_e^\pi\geq U$ at every active starting node and date implies that this transfer compensates the policy's loss everywhere in the declared nodal economy. We search for a common $e$ using monotonicity of CRRA felicity in consumption and verify the final inequality by outward propagation. Table~\ref{tab:r9welfare} gives the resulting upper percentages. These conservative transfers can be economically substantial even when the absolute utility bound is below $0.1$; the threshold is not described as a universally small welfare loss.
\begin{table}[htbp]
\centering\small\setlength{\tabcolsep}{4pt}
\caption{Sufficient financed flow-consumption supplements}\label{tab:r9welfare}
\begin{tabular}{lrrr}
\toprule
Proposal & Seed & Supplement, percent & Verification seconds \\
\midrule
Actor & 11 & 20.93 & 27.42 \\
Actor & 29 & 19.17 & 26.86 \\
Actor & 47 & 18.87 & 26.94 \\
Direct & 11 & 15.33 & 27.06 \\
Direct & 29 & 18.06 & 28.34 \\
Direct & 47 & 15.33 & 27.50 \\
Search & 11 & 20.44 & 31.32 \\
Search & 29 & 22.46 & 31.35 \\
Search & 47 & 19.42 & 31.37 \\
\bottomrule
\end{tabular}
\par\vspace{3pt}\begin{minipage}{0.97\linewidth}\footnotesize Uniform sufficient compensation over every active starting node and date. State dynamics, the original policy, terminal bequest, and adjustment cost stay fixed. Consumption is externally financed. These are not feasible self-financed policy improvements or common homothetic consumption equivalents.\end{minipage}
\end{table}

\subsection{A global benchmark for the coupled capital economy}\label{sec:r9capital}
A comparison of two frozen policies does not indicate how far either is from the economic optimum. We supplement the paired simulations by a bound on the original, continuous-time, dense nonlinear capital economy. It does not rely on a Riccati equation, separability of the optimal policy, a state mesh, or a sampled residual.

In the notation of the capital experiment, let $Y_i$ be log capital and $m_i\in[\underline m,\overline m]$ the consumption ratio. The flow payoff is $\overline{\log m}+\overline Y-(\zeta/2)\overline m^{\,2}$, and the terminal payoff is $\overline Y-\chi\overline{(Y-\overline Y)^2}$, with $\chi=0.03$. The drift of $Y_i$ is $a+\kappa\tanh((BY)_i)-m_i-(\sigma_I^2+\sigma_C^2)/2$. Common noise affects the mean but cancels from cross-sectional dispersion. Define
\[
 b_\pm=a\pm\kappa-\tfrac12(\sigma_I^2+\sigma_C^2),\qquad
 w(t)=\frac{1-e^{-\rho(T-t)}}{\rho}+e^{-\rho(T-t)},
\]
and let
\[
 m_0(t)=\operatorname{clip}_{[\underline m,\overline m]}
 \frac{2}{w(t)+\sqrt{w(t)^2+4\zeta}}.
\]
This common, deterministic schedule is feasible in the unchanged economy.

\begin{proposition}[Continuous-economy performance benchmark]\label{prop:r9capital}
Starting from $Y_i(0)=0$ for every $i$, let $V^*$ be the optimal expected payoff of the original $d$-dimensional capital economy. Define
\[
 U=\int_0^T e^{-\rho t}\{\log m_0(t)-\tfrac{\zeta}{2}m_0(t)^2+w(t)[b_+-m_0(t)]\}\,dt,
\]
\[
 D_d=\{\kappa T+\sigma_I\sqrt{(1-1/d)T}\}^{2},
\]
\[
 L_d=\int_0^T e^{-\rho t}\{\log m_0(t)-\tfrac{\zeta}{2}m_0(t)^2+w(t)[b_--m_0(t)]\}\,dt
 -\chi e^{-\rho T}D_d.
\]
Then $L_d\leq J(m_0)\leq V^*\leq U$ for every admissible dense coupling matrix $B$ in the stated model. The inequalities apply to the continuous economy on $\mathbb R^d$, not merely to an Euler simulation.
\end{proposition}

The bound exploits bounded nonlinear drift, concavity of consumption utility, and an explicit feasible policy. Its numerical integration uses outward rectangle sums, interval exponential and logarithm evaluations, and square roots checked by interval squaring. Table~\ref{tab:r9capital} reports the computed brackets. Their width is economically material and is not presented as a small regret bound for a learned policy. They nevertheless provide an absolute benchmark absent from a paired-policy comparison. Simulations of $m_0$ use the same Brownian paths as the saved neural policies. The resulting differences and confidence intervals concern their common Euler discretization; an unknown simulation bias is not subtracted from the analytical bound to create a neural-policy certificate.
\begin{table}[htbp]
\centering\small\setlength{\tabcolsep}{4pt}
\caption{Global bounds in the original continuous capital economy}\label{tab:r9capital}
\begin{tabular}{lrrrr}
\toprule
Dimension & Time panels & Feasible lower & Optimal upper & Width \\
\midrule
10 & 4096 & -1.489064 & -1.051803 & 0.437261 \\
10 & 16384 & -1.488883 & -1.051998 & 0.436884 \\
20 & 4096 & -1.489130 & -1.051803 & 0.437327 \\
20 & 16384 & -1.488949 & -1.051998 & 0.436950 \\
\bottomrule
\end{tabular}
\par\vspace{3pt}\begin{minipage}{0.97\linewidth}\footnotesize Initial log capital is zero in every coordinate; horizon one. The bounds cover every admissible control on $\mathbb R^d$. They are not small loss bounds for neural policies. The same feasible schedule is evaluated on the saved Brownian paths in the supplement.\end{minipage}
\end{table}



